\documentclass[../../main.tex]{subfiles}
\graphicspath{{\subfix{../../image/}}}

\begin{document}

\section{特征函数}\label{section:ch04-特征函数}

\subsection*{一、定义}

数字特征只反映了概率分布的某些侧面，一般并不能通过它们来完全确定分布函数，本节将要引进的特征函数，既能完全决定分布函数而又具有良好的分析性质。

为了定义特征函数，我们需要稍为拓广一下随机变量的概念，引进复随机变量。

\begin{definition}[复随机变量]\label{definition:ch04-复随机变量}
如果 $\xi$ 与 $\eta$ 都是概率空间 $(\Omega,\mathscr{F},P)$ 上的实值随机变量，则称 $\zeta=\xi+\mathrm{i}\eta$ 为\textbf{复随机变量}。
\end{definition}

从定义知道，对复随机变量的研究本质上是对二维随机向量的研究。这里举一个例子：如果二维向量 $(\xi_1,\eta_1)$ 与 $(\xi_2,\eta_2)$ 是独立的，则我们称复随机变量 $\zeta_1=\xi_1+\mathrm{i}\eta_1$ 与 $\zeta_2=\xi_2+\mathrm{i}\eta_2$ 是独立的。

定义一个复随机变量 $\zeta=\xi+\mathrm{i}\eta$ 的数学期望为
\[
E\zeta=E\xi+\mathrm{i}E\eta.
\]

对复随机变量也可以平行于实随机变量建立起一系列结果。例如，若 $\zeta_1,\zeta_2,\cdots,\zeta_n$ 是相互独立的，则
\[
E\zeta_1\zeta_2\cdots\zeta_n=E\zeta_1 E\zeta_2\cdots E\zeta_n.
\]

又如，若 $g(x)$ 是一个一元 Borel 可测函数，而 $\eta=g(\xi)$，则成立复佚名统计学家公式
\begin{align}
E\mathrm{e}^{\mathrm{i}t\eta}=E\mathrm{e}^{\mathrm{i}tg(\xi)}=\int_{-\infty}^{\infty}\mathrm{e}^{\mathrm{i}tg(x)}\,\mathrm{d}F_\xi(x).\label{align:ch04-复佚名公式}
\end{align}
这里使用 Euler 公式 $\mathrm{e}^{\mathrm{i}\theta}=\cos\theta+\mathrm{i}\sin\theta$。

以后将随时引用这类结果而不再加以说明。

下面引进随机变量 $\xi$ 的特征函数。

\begin{definition}[特征函数]\label{definition:ch04-特征函数}
若随机变量 $\xi$ 的分布函数为 $F_\xi(x)$，则称
\begin{align}
f_\xi(t)=E\mathrm{e}^{\mathrm{i}t\xi}=\int_{-\infty}^{\infty}\mathrm{e}^{\mathrm{i}tx}\,\mathrm{d}F_\xi(x)\label{align:ch04-特征函数定义}
\end{align}
为 $\xi$ 的\textbf{特征函数}（characteristic function）。
\end{definition}

特征函数是一个实变量的复值函数，由于 $|\mathrm{e}^{\mathrm{i}tx}|=1$，所以它对一切实数 $t$ 都有意义。

显然特征函数只与分布函数有关，因此亦称某一分布函数的特征函数。

对于离散型随机变量，若其分布列为
\begin{center}
\begin{tabular}{c|ccc}
$\xi$ & $x_1$ & $x_2$ & $\cdots$\\
\hline
$P$ & $p_1$ & $p_2$ & $\cdots$
\end{tabular}
\end{center}
则其特征函数为
\begin{align}
f(t)=\sum_{j=1}^{\infty}p_j\mathrm{e}^{\mathrm{i}t x_j}.\label{align:ch04-离散特征函数}
\end{align}
特别地，对于整值随机变量，若其母函数为 $P(s)$，则 $f(t)=P(\mathrm{e}^{\mathrm{i}t})$。

对于连续型随机变量，若其分布密度函数为 $p(x)$，则其特征函数为
\begin{align}
f(t)=\int_{-\infty}^{\infty}\mathrm{e}^{\mathrm{i}tx}p(x)\,\mathrm{d}x.\label{align:ch04-连续特征函数}
\end{align}
这时，特征函数是密度函数 $p(x)$ 的 Fourier（傅里叶）变换。

一般情况下的特征函数可以看作是这种 Fourier 变换的推广。Fourier 分析是数学中一种非常有力的工具，它在许多数学分支中都起了重大作用，以后我们将会看到，它在概率论中也占有突出的地位。

下面指出一些重要分布的特征函数。

\begin{example}[退化分布]\label{example:ch04-退化特征函数}
退化分布 $I_c(x)$ 的特征函数为
\begin{align}
f(t)=\mathrm{e}^{\mathrm{i}ct}.\label{align:ch04-退化特征函数}
\end{align}
\end{example}

\begin{example}[二项分布]\label{example:ch04-二项特征函数}
二项分布 $B(n,p)$ 的特征函数为
\begin{align}
f(t)=(p\mathrm{e}^{\mathrm{i}t}+q)^n.\label{align:ch04-二项特征函数}
\end{align}
\end{example}

\begin{example}[Poisson 分布]\label{example:ch04-Poisson特征函数}
Poisson 分布 $P(\lambda)$ 的特征函数为
\begin{align}
f(t)=\mathrm{e}^{\lambda(\mathrm{e}^{\mathrm{i}t}-1)}.\label{align:ch04-Poisson特征函数}
\end{align}
\end{example}

\begin{example}[$\Gamma$ 分布]\label{example:ch04-Gamma特征函数}
$\Gamma$ 分布 $\Gamma(r,\lambda)$ 的特征函数为
\begin{align}
f(t)&=\int_0^{\infty}\mathrm{e}^{\mathrm{i}tx}\frac{\lambda^r}{\Gamma(r)}x^{r-1}\mathrm{e}^{-\lambda x}\,\mathrm{d}x\nonumber\\
&=\int_0^{\infty}\frac{\lambda^r}{\Gamma(r)}x^{r-1}\mathrm{e}^{-\lambda(1-\frac{\mathrm{i}t}{\lambda})x}\,\mathrm{d}x\nonumber\\
&=\left(1-\frac{\mathrm{i}t}{\lambda}\right)^{-r}.\label{align:ch04-Gamma特征函数}
\end{align}

特别地，参数为 $\lambda$ 的指数分布 $\mathrm{Exp}(\lambda)$，即 $\Gamma(1,\lambda)$ 的特征函数为
\begin{align}
f(t)=\left(1-\frac{\mathrm{i}t}{\lambda}\right)^{-1}.\label{align:ch04-指数特征函数}
\end{align}

同样地，参数 $n$ 的 $\chi^2$ 分布，即 $\Gamma\left(\dfrac{n}{2},\dfrac12\right)$ 的特征函数为
\begin{align}
f(t)=(1-2\mathrm{i}t)^{-\frac{n}{2}}.\label{align:ch04-卡方特征函数}
\end{align}
\end{example}

\subsection*{二、性质}

下面讨论特征函数的一些基本性质。

\noindent\textbf{性质 1}\quad 特征函数 $f(t)$ 有如下性质：
\begin{align}
f(0)=1,\label{align:ch04-f0=1}\\
|f(t)|\leqslant f(0),\label{align:ch04-模长界}\\
f(-t)=\overline{f(t)}.\label{align:ch04-共轭}
\end{align}

\begin{proof}
\[
f(0)=\int_{-\infty}^{\infty}1\,\mathrm{d}F(x)=1,
\]
\[
|f(t)|\leqslant\int_{-\infty}^{\infty}|\mathrm{e}^{\mathrm{i}tx}|\,\mathrm{d}F(x)=1=f(0),
\]
\[
f(-t)=\int_{-\infty}^{\infty}\mathrm{e}^{-\mathrm{i}tx}\,\mathrm{d}F(x)
=\overline{\int_{-\infty}^{\infty}\mathrm{e}^{\mathrm{i}tx}\,\mathrm{d}F(x)}=\overline{f(t)}.
\]
\end{proof}

\noindent\textbf{性质 2}\quad 特征函数在 $(-\infty,\infty)$ 上一致连续。

\noindent\textbf{性质 3}\quad 对于任意的正整数 $n$ 及任意实数 $t_1,t_2,\cdots,t_n$ 及复数 $\lambda_1,\lambda_2,\cdots,\lambda_n$，成立
\begin{align}
\sum_{k=1}^n\sum_{j=1}^n f(t_k-t_j)\lambda_k\overline{\lambda_j}\geqslant 0.\label{align:ch04-非负定}
\end{align}

\begin{proof}
\begin{align*}
\sum_{k=1}^n\sum_{j=1}^n f(t_k-t_j)\lambda_k\overline{\lambda_j}
&=\sum_{k=1}^n\sum_{j=1}^n E\mathrm{e}^{\mathrm{i}(t_k-t_j)\xi}\lambda_k\overline{\lambda_j}\\
&=E\left\{\sum_{k=1}^n\sum_{j=1}^n\mathrm{e}^{\mathrm{i}(t_k-t_j)\xi}\lambda_k\overline{\lambda_j}\right\}\\
&=E\left\{\left(\sum_{k=1}^n\mathrm{e}^{\mathrm{i}t_k\xi}\lambda_k\right)\left(\sum_{j=1}^n\mathrm{e}^{-\mathrm{i}t_j\xi}\overline{\lambda_j}\right)\right\}\\
&=E\left|\sum_{k=1}^n\mathrm{e}^{\mathrm{i}t_k\xi}\lambda_k\right|^2\geqslant 0.
\end{align*}
\end{proof}

这个性质称为特征函数的\textbf{非负定性}，以后我们将会看到，这是特征函数最本质的性质之一。

\noindent\textbf{性质 4}\quad 两个相互独立的随机变量之和的特征函数等于它们的特征函数之积。

\begin{proof}
设 $\xi_1$ 与 $\xi_2$ 是两个相互独立的随机变量，而 $\eta=\xi_1+\xi_2$，由 $\xi_1$ 与 $\xi_2$ 的独立性不难推得复随机变量 $\mathrm{e}^{\mathrm{i}t\xi_1}$ 与 $\mathrm{e}^{\mathrm{i}t\xi_2}$ 也是独立的，因此
\[
E\mathrm{e}^{\mathrm{i}t\eta}=E\mathrm{e}^{\mathrm{i}t(\xi_1+\xi_2)}=E\mathrm{e}^{\mathrm{i}t\xi_1}\cdot E\mathrm{e}^{\mathrm{i}t\xi_2}.
\]
\end{proof}

性质 4 可用归纳法推广到 $n$ 个独立随机变量之和的场合。

应当着重指出，正是由于性质 4，才使特征函数在概率论中占有重要地位。由于这个性质，独立随机变量和的特征函数可以方便地用各个特征函数相乘来求得，而独立和的分布密度要通过卷积这种复杂的运算才能得到，相比之下，用特征函数来处理独立和问题就有力得多。独立和问题在概率论的古典问题中占有``中心''地位，而这些问题的解决大大有赖于特征函数的引进与使用。

\noindent\textbf{性质 5}\quad 设随机变量 $\xi$ 有 $n$ 阶矩存在，则它的特征函数可微分 $n$ 次，且当 $k\leqslant n$ 时：
\begin{align}
f^{(k)}(0)=\mathrm{i}^k E\xi^k.\label{align:ch04-特征函数导数}
\end{align}

\begin{proof}
\[
\left|\frac{\mathrm{d}^k}{\mathrm{d}t^k}\bigl(\mathrm{e}^{\mathrm{i}tx}\bigr)\right|
=\left|\mathrm{i}^k x^k\mathrm{e}^{\mathrm{i}tx}\right|\leqslant|x|^k.
\]
由于 $\xi$ 的 $k$ 阶矩存在，故 $\int_{-\infty}^{\infty}|x|^k\,\mathrm{d}F(x)<\infty$，因而可作下列积分号下的微分
\begin{align*}
f^{(k)}(t)&=\int_{-\infty}^{\infty}\frac{\mathrm{d}^k}{\mathrm{d}t^k}\bigl(\mathrm{e}^{\mathrm{i}tx}\bigr)\,\mathrm{d}F(x)\\
&=\mathrm{i}^k\int_{-\infty}^{\infty}x^k\mathrm{e}^{\mathrm{i}tx}\,\mathrm{d}F(x).
\end{align*}
取 $t=0$ 即得~\eqref{align:ch04-特征函数导数}。
\end{proof}

性质 5 使我们可以方便地求得随机变量的各阶矩。

\noindent\textbf{推论}\quad 若随机变量 $\xi$ 有 $n$ 阶矩存在，则它的特征函数可作如下展开：
\begin{align}
f(t)=1+(\mathrm{i}t)E\xi+\frac{(\mathrm{i}t)^2}{2!}E\xi^2+\cdots+\frac{(\mathrm{i}t)^n}{n!}E\xi^n+o(t^n).\label{align:ch04-特征函数泰勒}
\end{align}

\begin{proof}
由性质 5，$f(t)$ 可以在 $t=0$ 近旁作 Taylor 展开，公式~\eqref{align:ch04-特征函数泰勒} 就是在带有 Peano 余项的展开式中，代入~\eqref{align:ch04-特征函数导数} 式而得到的。
\end{proof}

\noindent\textbf{性质 6}\quad 设 $\eta=a\xi+b$，这里 $a,b$ 为常数，则
\begin{align}
f_\eta(t)=\mathrm{e}^{\mathrm{i}bt}f_\xi(at).\label{align:ch04-线性变换特征函数}
\end{align}

\begin{proof}
\begin{align*}
f_\eta(t)=E\mathrm{e}^{\mathrm{i}t\eta}=E\mathrm{e}^{\mathrm{i}t(a\xi+b)}
=\mathrm{e}^{\mathrm{i}bt}E\mathrm{e}^{\mathrm{i}ta\xi}=\mathrm{e}^{\mathrm{i}bt}f_\xi(at).
\end{align*}
\end{proof}

\begin{example}[正态分布的特征函数]\label{example:ch04-正态特征函数}
先讨论 $N(0,1)$ 的场合：
\[
f(t)=\frac{1}{\sqrt{2\pi}}\int_{-\infty}^{\infty}\mathrm{e}^{\mathrm{i}tx}\mathrm{e}^{-x^2/2}\,\mathrm{d}x
=\frac{1}{\sqrt{2\pi}}\int_{-\infty}^{\infty}\cos tx\cdot\mathrm{e}^{-x^2/2}\,\mathrm{d}x.
\]
由于正态分布一阶矩存在，可对上式求导，得
\begin{align*}
f'(t)&=\frac{1}{\sqrt{2\pi}}\int_{-\infty}^{\infty}(-x)\sin tx\cdot\mathrm{e}^{-x^2/2}\,\mathrm{d}x
=\frac{1}{\sqrt{2\pi}}\int_{-\infty}^{\infty}\sin tx\,\mathrm{d}\mathrm{e}^{-x^2/2}\\
&=\left[\frac{1}{\sqrt{2\pi}}\sin tx\cdot\mathrm{e}^{-x^2/2}\right]_{-\infty}^{\infty}
-\frac{1}{\sqrt{2\pi}}\int_{-\infty}^{\infty}t\cos tx\cdot\mathrm{e}^{-x^2/2}\,\mathrm{d}x\\
&=-tf(t).
\end{align*}
因此
\[
\ln f(t)=-\frac{t^2}{2}+c.
\]
由于 $f(0)=1$，所以 $c=0$。这样一来，
\begin{align}
f(t)=\mathrm{e}^{-t^2/2}.\label{align:ch04-标准正态特征函数}
\end{align}

一般 $N(\mu,\sigma^2)$ 的场合，利用性质 6 即得
\begin{align}
f(t)=\mathrm{e}^{\mathrm{i}\mu t-\frac12\sigma^2 t^2}.\label{align:ch04-一般正态特征函数}
\end{align}
\end{example}

\begin{theorem}[特征函数在 $\mathbb{R}$ 上一致连续]\label{theorem:ch04-特征函数一致连续}
设 $f(t)=\int_{-\infty}^{\infty}\mathrm{e}^{\mathrm{i}tx}p(x)\,\mathrm{d}x$ 为某概率分布的特征函数，则 $f(t)$ 在 $(-\infty,+\infty)$ 上一致连续。
\end{theorem}

\begin{proof}
对任意 $t,h\in\mathbb{R}$，
\begin{align*}
|f(t+h)-f(t)|
&=\left|\int_{-\infty}^{\infty}\bigl(\mathrm{e}^{\mathrm{i}(t+h)x}-\mathrm{e}^{\mathrm{i}tx}\bigr)\,\mathrm{d}F(x)\right|\\
&\leqslant\int_{-\infty}^{\infty}|\mathrm{e}^{\mathrm{i}hx}-1|\,\mathrm{d}F(x)\\
&\leqslant 2\int_{|x|\geqslant A}\,\mathrm{d}F(x)+\int_{-A}^{A}|\mathrm{e}^{\mathrm{i}hx}-1|\,\mathrm{d}F(x)\\
&=2\int_{|x|\geqslant A}\,\mathrm{d}F(x)+2\int_{-A}^{A}\left|\sin\frac{hx}{2}\right|\,\mathrm{d}F(x).
\end{align*}
上式右边已与 $t$ 无关；可选足够大的 $A$ 使 $\int_{|x|\geqslant A}\,\mathrm{d}F(x)$ 任意小，对选定的 $A$ 再选充分小的 $|h|$ 可使第二个积分也任意小，从而证明了结论。
\end{proof}

\subsection*{三、逆转公式与唯一性定理}

现在来证明特征函数与分布函数是相互唯一确定的，由分布函数决定特征函数是显然的，剩下来的是需要证明可由特征函数唯一决定分布函数。

下面定理的证明要用到如下数学分析的引理。

\begin{lemma}[Dirichlet 积分]\label{lemma:ch04-Dirichlet积分}
设 $x_1<x_2$，
\begin{align}
g(T,x,x_1,x_2)=\frac{1}{\pi}\int_0^T\left[\frac{\sin t(x-x_1)}{t}-\frac{\sin t(x-x_2)}{t}\right]\mathrm{d}t,\label{align:ch04-g函数}
\end{align}
则
\begin{align}
\lim_{T\to\infty}g(T,x,x_1,x_2)=
\begin{cases}
0,&x<x_1\text{ 或 }x>x_2,\\
\dfrac12,&x=x_1\text{ 或 }x=x_2,\\
1,&x_1<x<x_2.
\end{cases}\label{align:ch04-g极限}
\end{align}
\end{lemma}

\begin{proof}
从数学分析中知道 Dirichlet 积分
\begin{align}
D(\alpha)=\frac{1}{\pi}\int_0^{\infty}\frac{\sin\alpha t}{t}\,\mathrm{d}t=
\begin{cases}
\dfrac12,&\alpha>0,\\
0,&\alpha=0,\\
-\dfrac12,&\alpha<0.
\end{cases}\label{align:ch04-Dirichlet}
\end{align}
而
\[
\lim_{T\to\infty}g(T,x,x_1,x_2)=D(x-x_1)-D(x-x_2).
\]
分别考察 $x$ 在区间 $(x_1,x_2)$ 的端点及内外时相应 Dirichlet 积分的值即得~\eqref{align:ch04-g极限}。
\end{proof}

\begin{theorem}[逆转公式]\label{theorem:ch04-逆转公式}
设分布函数 $F(x)$ 的特征函数为 $f(t)$，又 $x_1,x_2$ 是 $F(x)$ 的连续点，则
\begin{align}
F(x_2)-F(x_1)=\lim_{T\to\infty}\frac{1}{2\pi}\int_{-T}^T\frac{\mathrm{e}^{-\mathrm{i}tx_1}-\mathrm{e}^{-\mathrm{i}tx_2}}{\mathrm{i}t}f(t)\,\mathrm{d}t.\label{align:ch04-逆转公式}
\end{align}
\end{theorem}

\begin{proof}
不妨设 $x_1<x_2$，记
\begin{align*}
I_T&=\frac{1}{2\pi}\int_{-T}^T\frac{\mathrm{e}^{-\mathrm{i}tx_1}-\mathrm{e}^{-\mathrm{i}tx_2}}{\mathrm{i}t}f(t)\,\mathrm{d}t\\
&=\frac{1}{2\pi}\int_{-T}^T\int_{-\infty}^{\infty}\frac{\mathrm{e}^{-\mathrm{i}tx_1}-\mathrm{e}^{-\mathrm{i}tx_2}}{\mathrm{i}t}\mathrm{e}^{\mathrm{i}tx}\,\mathrm{d}F(x)\,\mathrm{d}t.
\end{align*}

为证被积函数的有界性，用到不等式
\[
|\mathrm{e}^{\mathrm{i}\alpha}-1|\leqslant|\alpha|.
\]
事实上，对 $\alpha>0$
\[
|\mathrm{e}^{\mathrm{i}\alpha}-1|=\left|\int_0^\alpha\mathrm{e}^{\mathrm{i}x}\,\mathrm{d}x\right|\leqslant\int_0^\alpha|\mathrm{e}^{\mathrm{i}x}|\,\mathrm{d}x=\alpha.
\]
对 $\alpha\leqslant 0$，取共轭即知不等式也成立。

因此
\[
\left|\frac{\mathrm{e}^{-\mathrm{i}tx_1}-\mathrm{e}^{-\mathrm{i}tx_2}}{\mathrm{i}t}\mathrm{e}^{\mathrm{i}tx}\right|\leqslant x_2-x_1.
\]

交换上述二次积分顺序得到
\begin{align*}
I_T&=\frac{1}{2\pi}\int_{-\infty}^{\infty}\left[\int_{-T}^T\frac{\mathrm{e}^{-\mathrm{i}tx_1}-\mathrm{e}^{-\mathrm{i}tx_2}}{\mathrm{i}t}\mathrm{e}^{\mathrm{i}tx}\,\mathrm{d}t\right]\mathrm{d}F(x)\\
&=\frac{1}{2\pi}\int_{-\infty}^{\infty}\left[\int_0^T\frac{\mathrm{e}^{\mathrm{i}t(x-x_1)}-\mathrm{e}^{-\mathrm{i}t(x-x_1)}-\mathrm{e}^{\mathrm{i}t(x-x_2)}+\mathrm{e}^{-\mathrm{i}t(x-x_2)}}{\mathrm{i}t}\,\mathrm{d}t\right]\mathrm{d}F(x)\\
&=\frac1\pi\int_{-\infty}^{\infty}\left[\int_0^T\left(\frac{\sin t(x-x_1)}{t}-\frac{\sin t(x-x_2)}{t}\right)\mathrm{d}t\right]\mathrm{d}F(x)\\
&=\int_{-\infty}^{\infty}g(T,x,x_1,x_2)\,\mathrm{d}F(x).
\end{align*}

此处 $g(T,x,x_1,x_2)$ 按~\eqref{align:ch04-g函数} 定义。由~\eqref{align:ch04-g极限} 可以知道 $|g(T,x,x_1,x_2)|$ 有界，因此由 Lebesgue 控制收敛定理\footnote{为了避免冗长的分析论证，在证明特征函数的逆转公式及逆极限定理时，我们共四次使用了实变函数论中关于极限号与积分号交换的 Lebesgue 控制收敛定理。}并利用引理的结果可得：
\[
\lim_{T\to\infty}I_T=\int_{-\infty}^{\infty}\lim_{T\to\infty}g(T,x,x_1,x_2)\,\mathrm{d}F(x)
=F(x_2)-F(x_1).
\]
\end{proof}

\begin{theorem}[唯一性定理]\label{theorem:ch04-特征函数唯一性}
分布函数由其特征函数唯一决定。
\end{theorem}

\begin{proof}
应用逆转公式，在 $F(x)$ 的每一连续点上，当 $y$ 沿着 $F(x)$ 的连续点趋于 $-\infty$ 时，有
\begin{align}
F(x)=\frac{1}{2\pi}\lim_{y\to-\infty}\lim_{T\to\infty}\int_{-T}^T\frac{\mathrm{e}^{-\mathrm{i}ty}-\mathrm{e}^{-\mathrm{i}tx}}{\mathrm{i}t}f(t)\,\mathrm{d}t,\label{align:ch04-分布函数由特征函数}
\end{align}
而分布函数由其连续点上的值唯一决定。
\end{proof}

由唯一性定理可知特征函数也完整地描述了随机变量。

特别当 $f(t)$ 是绝对可积函数时，有下列更强的结果。

\begin{theorem}[密度的反演]\label{theorem:ch04-密度反演}
若 $\int_{-\infty}^{\infty}|f(t)|\,\mathrm{d}t<\infty$，则相应的分布函数 $F(x)$ 的导数存在并连续，而且
\begin{align}
F'(x)=\frac{1}{2\pi}\int_{-\infty}^{\infty}\mathrm{e}^{-\mathrm{i}tx}f(t)\,\mathrm{d}t.\label{align:ch04-密度反演公式}
\end{align}
\end{theorem}

\begin{proof}
由逆转公式，若 $x+\Delta x$ 及 $x$ 是 $F(x)$ 的连续点，则
\[
\frac{F(x+\Delta x)-F(x)}{\Delta x}
=\lim_{T\to\infty}\frac{1}{2\pi}\int_{-T}^T\frac{\mathrm{e}^{-\mathrm{i}tx}-\mathrm{e}^{-\mathrm{i}t(x+\Delta x)}}{\mathrm{i}t\Delta x}f(t)\,\mathrm{d}t.
\]
利用 $|\mathrm{e}^{\mathrm{i}\alpha}-1|\leqslant|\alpha|$，可得
\[
\left|\frac{\mathrm{e}^{-\mathrm{i}tx}-\mathrm{e}^{-\mathrm{i}t(x+\Delta x)}}{\mathrm{i}t\Delta x}\right|\leqslant 1.
\]

依假设 $\int_{-\infty}^{\infty}|f(t)|\,\mathrm{d}t<\infty$，因此
\[
\frac{F(x+\Delta x)-F(x)}{\Delta x}
=\frac{1}{2\pi}\int_{-\infty}^{\infty}\frac{\mathrm{e}^{-\mathrm{i}tx}-\mathrm{e}^{-\mathrm{i}t(x+\Delta x)}}{\mathrm{i}t\Delta x}f(t)\,\mathrm{d}t.
\]

利用控制收敛定理
\begin{align*}
F'(x)&=\lim_{\Delta x\to 0}\frac{F(x+\Delta x)-F(x)}{\Delta x}\\
&=\frac{1}{2\pi}\int_{-\infty}^{\infty}\lim_{\Delta x\to 0}\frac{\mathrm{e}^{-\mathrm{i}tx}-\mathrm{e}^{-\mathrm{i}t(x+\Delta x)}}{\mathrm{i}t\Delta x}f(t)\,\mathrm{d}t\\
&=\frac{1}{2\pi}\int_{-\infty}^{\infty}\mathrm{e}^{-\mathrm{i}tx}f(t)\,\mathrm{d}t.
\end{align*}
因此 $p(x)=F'(x)$ 存在而有界。再次利用控制收敛定理就可得 $F'(x)$ 的连续性。

因此在 $f(t)$ 是绝对可积的条件下，分布密度 $p(x)$ 与特征函数 $f(t)$ 通过 Fourier 变换来联系。
\end{proof}

\subsection*{四、分布函数的再生性}

许多重要的分布函数具有一个有趣的性质——\textbf{再生性}。这个性质用特征函数来研究最为方便。下面通过几个例子来说明它。

\begin{example}[二项分布的再生性]\label{example:ch04-二项再生}
若 $\xi_1$ 服从 $B(m,p)$，$\xi_2$ 服从 $B(n,p)$，而且 $\xi_1$ 与 $\xi_2$ 独立，则 $\eta=\xi_1+\xi_2$ 服从 $B(m+n,p)$。

事实上 $f_{\xi_1}(t)=(p\mathrm{e}^{\mathrm{i}t}+q)^m$，$f_{\xi_2}(t)=(p\mathrm{e}^{\mathrm{i}t}+q)^n$，由性质 4 知
\[
f_\eta(t)=(p\mathrm{e}^{\mathrm{i}t}+q)^{m+n}.
\]
因此由唯一性定理知 $\eta$ 服从 $B(m+n,p)$。

这个事实简记作
\begin{align}
B(n_1,p)*B(n_2,p)=B(n_1+n_2,p).\label{align:ch04-二项再生}
\end{align}
\end{example}

\begin{example}[Poisson 分布的再生性]\label{example:ch04-Poisson再生}
若 $\xi_1$ 服从 $P(\lambda_1)$，$\xi_2$ 服从 $P(\lambda_2)$，而且 $\xi_1$ 与 $\xi_2$ 独立，则 $\eta=\xi_1+\xi_2$ 服从 $P(\lambda_1+\lambda_2)$。

事实上
\[
f_{\xi_1}(t)=\mathrm{e}^{\lambda_1(\mathrm{e}^{\mathrm{i}t}-1)},\qquad
f_{\xi_2}(t)=\mathrm{e}^{\lambda_2(\mathrm{e}^{\mathrm{i}t}-1)},
\]
\[
f_\eta(t)=\mathrm{e}^{(\lambda_1+\lambda_2)(\mathrm{e}^{\mathrm{i}t}-1)}.
\]

这个事实简记作
\begin{align}
P(\lambda_1)*P(\lambda_2)=P(\lambda_1+\lambda_2).\label{align:ch04-Poisson再生}
\end{align}
\end{example}

\begin{example}[正态分布的再生性]\label{example:ch04-正态再生}
若 $\xi_1$ 服从 $N(\mu_1,\sigma_1^2)$，$\xi_2$ 服从 $N(\mu_2,\sigma_2^2)$，而且 $\xi_1$ 与 $\xi_2$ 独立，则 $\eta=\xi_1+\xi_2$ 服从 $N(\mu_1+\mu_2,\sigma_1^2+\sigma_2^2)$。

事实上
\[
f_{\xi_1}(t)=\mathrm{e}^{\mathrm{i}\mu_1 t-\frac12\sigma_1^2 t^2},\qquad
f_{\xi_2}(t)=\mathrm{e}^{\mathrm{i}\mu_2 t-\frac12\sigma_2^2 t^2},
\]
\[
f_\eta(t)=\mathrm{e}^{\mathrm{i}(\mu_1+\mu_2)t-\frac12(\sigma_1^2+\sigma_2^2)t^2}.
\]

这个事实简记作
\begin{align}
N(\mu_1,\sigma_1^2)*N(\mu_2,\sigma_2^2)=N(\mu_1+\mu_2,\sigma_1^2+\sigma_2^2).\label{align:ch04-正态再生}
\end{align}
\end{example}

\begin{example}[$\Gamma$ 分布的再生性]\label{example:ch04-Gamma再生}
若 $\xi_1$ 服从 $\Gamma(r_1,\lambda)$，$\xi_2$ 服从 $\Gamma(r_2,\lambda)$，而且 $\xi_1$ 与 $\xi_2$ 独立，则 $\eta=\xi_1+\xi_2$ 服从 $\Gamma(r_1+r_2,\lambda)$。

事实上，由~\eqref{align:ch04-Gamma特征函数}
\[
f_{\xi_1}(t)=\left(1-\frac{\mathrm{i}t}{\lambda}\right)^{-r_1},\qquad
f_{\xi_2}(t)=\left(1-\frac{\mathrm{i}t}{\lambda}\right)^{-r_2},
\]
\[
f_\eta(t)=\left(1-\frac{\mathrm{i}t}{\lambda}\right)^{-(r_1+r_2)}.
\]

这个事实简记作
\begin{align}
\Gamma(r_1,\lambda)*\Gamma(r_2,\lambda)=\Gamma(r_1+r_2,\lambda).\label{align:ch04-Gamma再生}
\end{align}

特别地，$\chi_n^2$ 分布即为 $\Gamma\left(\dfrac{n}{2},\dfrac12\right)$，也具有再生性：
\begin{align}
\chi_m^2*\chi_n^2=\chi_{m+n}^2.\label{align:ch04-卡方可加}
\end{align}
这个性质，通常称为 $\chi^2$ 分布的可加性，已见于式~(3.3.38)，在数理统计中是最常使用的命题之一。
\end{example}

还有不少重要分布也有再生性，留给读者作为练习。

还有人研究了这类命题的逆命题——分布函数的分解问题，即若两个独立随机变量之和服从某一分布，问是否能断定这两个随机变量也分别服从这个分布。已经证明对于正态分布及 Poisson 分布逆命题的确成立。

\subsection*{五、多元特征函数}

若随机向量 $(\xi_1,\xi_2,\cdots,\xi_n)$ 的分布函数为 $F(x_1,x_2,\cdots,x_n)$，与随机变量相仿，我们可以定义它的特征函数
\begin{align}
f(t_1,t_2,\cdots,t_n)=\int_{-\infty}^{\infty}\cdots\int_{-\infty}^{\infty}\mathrm{e}^{\mathrm{i}(t_1x_1+\cdots+t_nx_n)}\,\mathrm{d}F(x_1,\cdots,x_n).\label{align:ch04-多元特征函数定义}
\end{align}

可以类似于一元的场合，建立起 $n$ 元特征函数的理论，由于方法完全相同，我们只叙述一些有关结论，证明一概从略。

\noindent\textbf{性质 1}\quad $f(t_1,t_2,\cdots,t_n)$ 在 $\mathbb{R}^n$ 中一致连续，而且
\[
|f(t_1,t_2,\cdots,t_n)|\leqslant f(0,0,\cdots,0)=1,
\]
\[
f(-t_1,-t_2,\cdots,-t_n)=\overline{f(t_1,t_2,\cdots,t_n)}.
\]

\noindent\textbf{性质 2}\quad 如果 $f(t_1,t_2,\cdots,t_n)$ 是 $(\xi_1,\xi_2,\cdots,\xi_n)$ 的特征函数，则 $\eta=a_1\xi_1+a_2\xi_2+\cdots+a_n\xi_n$ 的特征函数为
\[
f_\eta(t)=f(a_1 t,a_2 t,\cdots,a_n t).
\]

\noindent\textbf{性质 3}\quad 如果矩 $E\xi_1^{k_1}\xi_2^{k_2}\cdots\xi_n^{k_n}$ 存在，则
\begin{align}
E\xi_1^{k_1}\xi_2^{k_2}\cdots\xi_n^{k_n}
=\mathrm{i}^{-\sum_{j=1}^n k_j}\left.\left[\frac{\partial^{k_1+k_2+\cdots+k_n}}{\partial t_1^{k_1}\partial t_2^{k_2}\cdots\partial t_n^{k_n}}f(t_1,t_2,\cdots,t_n)\right]\right|_{t_1=t_2=\cdots=t_n=0}.\label{align:ch04-多元矩}
\end{align}

\noindent\textbf{性质 4}\quad 若 $(\xi_1,\xi_2,\cdots,\xi_n)$ 的特征函数为 $f(t_1,t_2,\cdots,t_n)$，则 $k(k<n)$ 维随机向量 $(\xi_1,\xi_2,\cdots,\xi_k)$ 的特征函数为
\[
f_{1,2,\cdots,k}(t_1,t_2,\cdots,t_k)=f(t_1,t_2,\cdots,t_k,0,\cdots,0).
\]
这是前 $k$ 个分量的 $k$ 元边际分布函数对应的特征函数。对应于任意 $k$ 个分量 $\xi_{j_1},\xi_{j_2},\cdots,\xi_{j_k}$ 的边际分布函数的特征函数，可以类似得到。

\noindent\textbf{逆转公式}\quad 如果 $f(t_1,t_2,\cdots,t_n)$ 是随机向量 $(\xi_1,\xi_2,\cdots,\xi_n)$ 的特征函数，而 $F(x_1,x_2,\cdots,x_n)$ 是它的分布函数，则
\begin{align*}
P\{a_k\leqslant\xi_k<b_k,\ k=1,2,\cdots,n\}
&=\lim_{\substack{T_j\to\infty\\j=1,\cdots,n}}\frac{1}{(2\pi)^n}\int_{-T_1}^{T_1}\int_{-T_2}^{T_2}\cdots\int_{-T_n}^{T_n}\prod_{k=1}^n\frac{\mathrm{e}^{-\mathrm{i}t_k a_k}-\mathrm{e}^{-\mathrm{i}t_k b_k}}{\mathrm{i}t_k}\\
&\quad\cdot f(t_1,t_2,\cdots,t_n)\,\mathrm{d}t_1\mathrm{d}t_2\cdots\mathrm{d}t_n,
\end{align*}
其中 $a_k$ 和 $b_k$ 都是任意实数，但满足唯一的要求：$(\xi_1,\xi_2,\cdots,\xi_n)$ 落在平行体
\[
a_k\leqslant x_k<b_k,\qquad k=1,2,\cdots,n,
\]
的面上的概率等于零。

\noindent\textbf{唯一性定理}\quad 分布函数 $F(x_1,x_2,\cdots,x_n)$ 由其特征函数唯一决定。

有了唯一性定理，可以进一步证明特征函数的如下两个性质，它们表征了独立性。

\noindent\textbf{性质 5}\quad 若 $(\xi_1,\xi_2,\cdots,\xi_n)$ 的特征函数为 $f(t_1,t_2,\cdots,t_n)$，而 $\xi_j$ 的特征函数为 $f_{\xi_j}(t)$，$j=1,2,\cdots,n$，则随机变量 $\xi_1,\xi_2,\cdots,\xi_n$ 相互独立的充要条件为
\begin{align}
f(t_1,t_2,\cdots,t_n)=f_{\xi_1}(t_1)f_{\xi_2}(t_2)\cdots f_{\xi_n}(t_n).\label{align:ch04-独立乘积特征函数}
\end{align}

\noindent\textbf{性质 6}\quad 若以 $f_1(t_1,\cdots,t_n)$，$f_2(u_1,\cdots,u_m)$ 及 $f(t_1,\cdots,t_n,u_1,\cdots,u_m)$ 分别记随机向量 $(\xi_1,\cdots,\xi_n)$，$(\eta_1,\cdots,\eta_m)$ 及 $(\xi_1,\cdots,\xi_n,\eta_1,\cdots,\eta_m)$ 的特征函数，则 $(\xi_1,\cdots,\xi_n)$ 与 $(\eta_1,\cdots,\eta_m)$ 独立的充要条件为：对一切实数 $t_1,\cdots,t_n$ 及 $u_1,\cdots,u_m$ 成立
\begin{align}
f(t_1,\cdots,t_n,u_1,\cdots,u_m)=f_1(t_1,\cdots,t_n)f_2(u_1,\cdots,u_m).\label{align:ch04-随机向量独立特征函数}
\end{align}

在下一节，我们要用到如下定理，相应于此定理的一维结果，将在下章叙述并证明。

\begin{theorem}[连续性定理]\label{theorem:ch04-连续性定理}
若特征函数列 $\{f_k(t_1,t_2,\cdots,t_n)\}$ 收敛于一个连续函数 $f(t_1,t_2,\cdots,t_n)$，则函数 $f(t_1,t_2,\cdots,t_n)$ 是某分布函数所对应的特征函数。
\end{theorem}

\end{document}